\chapter{Attending to Remember: Recent Advances in Methods and Theory}

Hallmarks of human cognition, such as forming and carrying out complex plans in pursuit of goals or flexibly engaging in intricate and dynamic social interactions, are supported in part by an ensemble of neurocognitive mechanisms that enable episodic memory—that is, the ability to learn and later bring back to mind details from past experiences. Although much progress has been made in understanding episodic memory, there remain open questions about some of the key mechanisms that impact remembering or forgetting in any given moment. Attention, in particular, is one set of processes implicated in learning and remembering since the earliest investigations into episodic memory (Ebbinghaus, 1885/2013). Over the ensuing decades, remarkable advances have been made in understanding multiple forms of attention and how attention mechanisms modulate learning and lead to variability in whether and how we remember. Here, we highlight recent methodological advances and concomitant theoretical insights into how attention impacts learning and remembering and note important open questions and promising future directions.

\section{Interactions Between Neural Networks of Attention, Goal-State Processes, and Memory}

Cognitive scientists and neuroscientists have made tremendous progress in specifying and measuring different forms of attention, their neural mechanisms, as well as their interactions with related cognitive control processes. A central insight is the dichotomous nature of attention. Namely, there is top-down (goal-directed or endogenous) orienting of attention to and selection of goal-relevant stimuli, which contrasts with bottom-up (stimulus-driven or exogenous) attentional capture (e.g., Corbetta \& Shulman, 2002; Posner \& Petersen, 1990). To illustrate, consider this real-world driving scenario: While preparing to turn onto a specific street, top-down selective attention is directed toward street-name signs, whereas bottom-up attentional capture is prompted by the unexpected event of a dog running out from between two parked cars. Top-down and bottom-up attention systems dynamically interact with related mechanisms of cognitive control, including those that subserve the representation of goals and enable goals to govern attention and information processing. The interactions between attention and cognitive control are bidirectional: Not only do goal states influence attention, but attention also impacts goal representation.

As depicted in Figure 1a, there are three frontoparietal networks central to attention and cognitive control: top-down attention via the dorsal attention network, bottom-up attention via the ventral attention network, and the cognitive control network (also known as the frontoparietal control network; Cole \& Schneider, 2007; Corbetta \& Shulman, 2002; Corbetta et al., 2008; Menon \& D’Esposito, 2022). With respect to episodic memory, attention and cognitive control mechanisms can affect the representations of perceived and retrieved event features in the neocortex along with memory-relevant computations and representations within the medial temporal lobe (e.g., Cabeza et al., 2008; Dobbins \& Wagner, 2005; Hutchinson et al., 2014; Uncapher \& Wagner, 2009). Both the intensity of attention as well as its selectivity are at least two ways in which attention impacts memory encoding and retrieval.

New insights into interactions between attention, cognitive control, and memory have come from studies leveraging readouts of attention and/or goal states during the acquisition and expression of episodic memories. This includes utilizing temporally resolved psychophysiological tools—such as reaction-time variability, pupil diameter, and posterior alpha (8–12 Hz) power (see Table 1) assayed via scalp EEG—to measure moment-to-moment attentional fluctuations and between-individuals attentional variability. Adopting these tools, multiple studies have revealed that the strength of top-down attention just prior to and during learning or an attempt to remember correlates with memory performance (i.e., readiness to learn and readiness to remember; e.g., Cohen et al., 2015; Madore et al., 2020; Madore \& Wagner, 2022; Miller \& Unsworth, 2020; Miller et al., 2019; Robison et al., 2022).

To illustrate, one recent experiment examined interactions between attention, goal coding, and memory using a goal-directed associative memory task in which during each retrieval trial (Fig. 1b) participants were asked to indicate whether they remembered a test probe as having been previously encountered in one of two task contexts during an immediately preceding study phase. \footnote{During the memory encoding (study) phase, participants encountered images of everyday objects that were physically small (e.g., fixed at 150 × 150 pixels) or large (e.g., fixed at 450 × 450 pixels) on the screen on any given trial. Some of the objects were normatively more pleasant than others. On half of the study trials, participants were instructed to decide whether the upcoming object would be physically small or large; on the other half of trials, participants were instructed to indicate whether the upcoming object would be pleasant or unpleasant.} During the retrieval phase, participants were cued with one of three retrieval goals on any given trial. Readouts of top-down attention just prior to goal cueing included pupil size and EEG posterior alpha power, and the strength of goal coding was measured via a midfrontal event-related potential elicited by the retrieval goal cue.\footnote{An alternative approach to measuring the strength of goal representations is to use pattern-classification methods to quantify the strength of goal representation using neural-activity patterns within the frontoparietal control network (e.g., Waskom et al., 2014).} Analyses revealed that goal-coding strength varied as a function of the level of attention evident in the moment just prior to retrieval goal onset and that these interactions between attention and goal coding predicted whether retrieval would be successful or unsuccessful (Fig. 1c; Madore \& Wagner, 2022; Madore et al., 2020). Thus, attention impacts retrieval success in part by affecting the representation and maintenance of one’s mnemonic goal.

Other work utilizing temporally resolved psychophysiological tools has revealed additional ways in which memory processes regulate attentional control, such as when mnemonic prediction errors signal stimulus or event salience, leading to attention orienting (den Ouden et al., 2012) and, in turn, influencing the encoding of unexpected information (Bein et al., 2021). Complementing these findings, new data also indicate that the top-down/bottom-up dichotomy of attentional control is insufficient for explaining situations for which neither goals nor salience account for biases in selective attention (e.g., rewards associated with equally salient stimuli in conflict with current selection goals; Awh et al., 2012). Selection history (broadly construed) is argued to be a missing construct of attentional control (for a review, see Anderson et al., 2021), in which mnemonic traces of prior experience (across varying timescales) enable memory-guided attention that can not only reconstitute (i.e., “re-member”) relevant representations of past experiences but also leverage forward-looking memory traces (i.e., “pre-member”) that functionally interact with incoming sensory signals to prospectively regulate attentional control and guide behavior (Hutchinson \& Turk-Browne, 2012; Nobre \& Stokes, 2019). Integrating prediction error and selection history accounts with computational formalism promises to yield more comprehensive theories of the dynamic interactions of attention and episodic memory and their consequences for learning and remembering.

\section{Rhythmic Nature of Attention and Memory}

Understanding how attention impacts memory is further enabled by characterizing their temporal dynamics. Growing evidence, primarily from temporally dense sampling of behavior, suggests that attention operates rhythmically predominantly in the theta (between approximately 4 and 7 Hz) and/or alpha (between approximately 8 and 12 Hz) frequency ranges (e.g., Fiebelkorn \& Kastner, 2019; VanRullen, 2016; VanRullen \& Dubois, 2011). Critically, a key assumption is that the cue resets the ongoing phase of the attention rhythm; therefore, by systematically varying the delay between the cue and the stimulus, one can sample different phases of the attention rhythm. Even under conditions in which attention is supposedly sustained, such as during trials with valid cues (meaning attention is directed to the location where the target ultimately appears), rhythmic attention fluctuation can still be observed. This implies that there are optimal and suboptimal phases of attention for information processing alternating in cycles of about 200 ms. Importantly, given interactions between attention and memory, this raises fundamental questions about potential rhythmicity in memory behavior (Biba et al., 2024; Ter Wal et al., 2021), its underlying neural mechanisms, and the impact of rhythms of attention on memory function. A key open question is whether memory encoding and retrieval depend, in part, on the phase of the ongoing attentional rhythm at which to-be-encoded information, retrieval cues, or retrieval products fall.

Recent data reveal rhythmicity in episodic memory behavior (Biba et al., 2024; Ter Wal et al., 2021), with findings interpreted in the context of a prominent model of hippocampal theta in which opposite phases of hippocampal theta are posited to be differentially optimal for encoding versus retrieval—the separate phases of encoding and retrieval (SPEAR) model (Fig. 2a; Hasselmo et al., 2002). From the SPEAR perspective, the relative influences of the monosynaptic and trisynaptic pathways of the hippocampus are thought to differ with phase, prioritizing either stimulus/event input from the entorhinal cortex in support of encoding or internally generated mnemonic predictions from hippocampal subfield CA3 for retrieval, respectively. The strength of inputs along the trisynaptic pathway (e.g., from the dentate gyrus and entorhinal cortex to CA3) is additionally thought to differ with theta phase, driving CA3 to either retrieve or encode information (Kunec et al., 2005). The theta-phase dependency of encoding and retrieval computations may not only account for rhythmicity in memory behavior (Biba et al., 2024; Ter Wal et al., 2021) but also relate to retrieval-driven versus novelty-driven eye movements (Kragel et al., 2020).

Importantly, although the hippocampal theta phase (Saint Amour di Chanaz et al., 2023) has been linked with encoding and retrieval success, direct evidence for the coupling of specific hippocampal phases with behavioral rhythmicity remains limited. Furthermore, there may exist encoding and retrieval modes with effects that persist over seconds (Duncan et al., 2012); how these prolonged mnemonic modes relate to subsecond theta-specific oscillations in hippocampal computations and behavior remains an open question. Given the nascent literature on rhythms of memory, such behavioral oscillations could be linked, at least in part, to rhythms in attention (Biba et al., 2024). Moreover, neural substrates of behavioral rhythms in memory may reside in the hippocampus and/or in frontoparietal attention networks.

Separately, an extensive amount of research demonstrates that theta power is linked to episodic memory functions (Hsieh \& Ranganath, 2014; Herweg et al., 2020), with scalp EEG showing increases in theta power during successful encoding and retrieval (Fig. 2b) and intracranial EEG data from the human hippocampus and neocortex showing similar effects during associative recall (Herweg et al., 2020; Maoz et al., 2023). Recent findings also document how the hippocampus and neocortex interact during encoding and retrieval (Fernandez et al., 2024; Theves et al., 2024; Zhao \& Kuhl, 2024), which invites questions regarding the impact of hippocampal theta oscillations on cortical representations (Hanslmayr et al., 2024). Emerging evidence from scalp EEG suggests that the strength of encoded and retrieved event content in the neocortex, read out by machine-learning pattern-classification methods (Table 1), oscillates at a theta frequency (Fig. 2c; Kerrén et al., 2018). Notably, peaks in cortical representational strength during learning and remembering may couple with opposing phases of a virtual hippocampal theta, although the relationship between scalp-recorded theta and hippocampal theta remains unclear (Herweg et al., 2020; Mitchell et al., 2008).

Given the uncertainty about whether oscillations of event-content representations in the neocortex relate to hippocampal theta, future research should explore whether fluctuations in the strength of to-be-encoded and retrieved cortical representations have a nonhippocampal source and are governed instead by the phase of theta oscillations that support attentional sampling (Busch \& VanRullen, 2010; Fiebelkorn \& Kastner, 2019). Moreover, future research can extend recent behavioral findings that attentional and mnemonic processes fluctuate at similar frequencies to examine their possible connections through an underlying neural rhythmicity. Recent technological advances permitting closed-loop intracranial EEG recordings during ambulatory navigation (Maoz et al., 2023) could be leveraged to address these open questions more directly (for more information on closed-loop approaches, see Testing the Causality of Attention for Remembering section).

\section{Memory Precision in Aging}
Aging is accompanied by changes in specific cognitive faculties, including attention and episodic memory (e.g., Fortenbaugh et al., 2015; Hedden \& Gabrieli, 2004). Representational quality—including event-based feature representation (Table 1) during the encoding of experiences and reinstatement of previously encoded representations during retrieval—is central in many theoretical accounts and empirical investigations of age-related changes in episodic memory (e.g., Stark et al., 2019; Theves et al., 2024; Trelle et al., 2020). Fortunately, advances in paradigm design and analytics promise increased sensitivity in the detection of subtle memory changes. A prominent example is the measurement of memory precision, which was originally developed to examine representational quality and quantity in working memory (e.g., Ma et al., 2014). Unlike classic memory tasks, in which participants indicate their memory by selecting from among a few discrete categorical decisions, assays of memory precision task participants with indicating their memory for fine-grained details using more continuous response options. For instance, during learning, participants might encounter common objects, each presented in a random color sampled from 360 possible colors; then, during retrieval, the precision of memory for an object’s study-phase color is probed by having participants indicate their memory by selecting a color using a 360-degree color wheel (Fig. 3; e.g., Sutterer \& Awh, 2016). This approach can be generalized to probe memory precision for any feature that can be mapped onto an approximately perceptually uniform space, such as location, orientation, or shape (Cooper \& Ritchey, 2019; Li et al., 2020; Richter et al., 2016).

Research on age-related episodic memory decline demonstrates the utility and promise of probing memory precision (e.g., Korkki et al., 2020; Nilakantan et al., 2018). For example, Korkki et al. (2020) tasked young and older participants to encode the location, color, and orientation of objects and probed subsequent memory precision for each of the three features. They fit the retrieval data with a mixture model (Zhang \& Luck, 2008), which some posit separates guesses from memory judgments varying in precision, and found that the precision estimate declined consistently with age across all three feature dimensions. Richter et al. (2016) leveraged a memory precision paradigm and mixture modeling to demonstrate that different expressions of episodic memory map onto distinct neural substrates, with functional MRI activity in the hippocampus showing a categorical effect (being more active when any memory information was retrieved regardless of its precision) and activity in the angular gyrus showing a continuous effect (scaling with the precision of retrieved mnemonic features). Korkki et al. (2023) replicated this relationship between episodic memory precision and angular gyrus activity and further observed that variability in precision may relate to neocortical structural variability (Table 1).

Notably, in Souza et al. (2024), older adults benefited more than young adults from a retro-cue3 directing their attention toward a particular to-be-remembered stimulus out of several others just prior to being probed about their memory for the stimulus (i.e., memory precision for the stimulus’ color). By contrast, when no retro-cue was provided (i.e., when participants did not know which of the several stimuli would be tested in the seconds after the initial stimulus array presentation), young adults outperformed older adults in terms of memory precision. Although the delay between encoding and retrieval in Souza et al. (2024) was much shorter than that in Korkki et al. (2023), given that multiple stimuli were presented within and across memory encoding trials prior to the retrieval trials in Korkki et al. (2023), this pattern of results raises the possibility that age-related declines in memory precision may be attributed, in part, to age-related differences in selective attention.

One theory-informing application of memory precision paradigms could be to shed new light on neural dedifferentiation and its links to attention in aging (e.g., Koen et al., 2020; Park \& McDonough, 2013). Neural dedifferentiation is a prominent age-related change in the neural underpinnings of perception and episodic memory. It consists of a reduction in the selectivity of cortical activity with age, is posited to decrease the fidelity (or precision) of memory representations, and may, in part, result from declines in attention. Indeed, recent findings indicate that, in cognitively unimpaired older adults, age-related declines in memory are explained, in part, by two independent pathways: a decline in memory-related encoding activity in the dorsal attention network that in turn accounts for declines in neural selectivity, and the presence of preclinical Alzheimer’s disease pathology (as evidenced by plasma pTau181) that separately accounts for declines in neural selectivity (Sheng et al., 2024). A test of this account of dedifferentiation could come from integrating memory precision assays with measures of neural representational precision and readouts of attention (e.g., from pupillometry and/or EEG); doing so promises to deepen understanding of the multiple drivers of neurocognitive aging.

\section{Testing the Causality of Attention for Remembering}
Historically, attempts to investigate the causal relationship between attention and memory have relied on manipulations in which a participant’s attention is divided across a primary and secondary task (e.g., Baddeley et al., 1984; Craik et al., 1996; Murdock, 1965) or in which two feature dimensions—one relevant and one irrelevant on any given trial—compete for attention (e.g., Uncapher \& Rugg, 2009). This experimental approach has yielded a rich literature revealing worse memory performance when attention is divided versus when fully focused on memory-relevant information.

In this article, much of our discussion has centered on correlational observations of attention-memory interactions. Such findings, although informative, do not permit causal inferences. Over the last few decades, advances in computing and closed-loop interfaces (Table 1) have enabled a complementary approach to bridging this gap (Fig. 4a; e.g., deBettencourt et al., 2015, 2018, 2019; Keene et al., 2022; Salari \& Rose, 2016; Yoo et al., 2012). To illustrate, deBettencourt et al. (2019) examined how attentional states impact working memory using a real-time adaptive approach. Through moment-to-moment sampling of response time on a sustained attention task, deBettencourt et al. (2019) tracked intrinsic fluctuations of attention, detecting when a participant was in exceptionally high- or low-attention states. In turn, they delivered working memory probes during these optimal and suboptimal attention states4 and observed superior performance in the former. Notably, this closed-loop approach generally affords greater control over factors of interest along with increased power.

Adaptive experimental approaches can be extended to causal investigations of episodic memory-attention interactions (Fig. 4a). Using a real-time framework, one (or multiple concurrent) psychophysiological readouts of fluctuations in sustained attention (e.g., via pupillometry; see Keene et al., 2022) could control either the temporal delivery of events to optimal versus suboptimal moments or trigger salient reorienting cues to reengage attention during momentary lapses just prior to engaging in the act of encoding or retrieval (Figs. 4b and 4c). However, despite the advantages afforded by closed-loop approaches for more precise causal investigations, unforeseen challenges may arise when designing such experiments. When developing a closed-loop pipeline, researchers should adhere to signal processing best practices/limitations, rigorous testing of computational/hardware integration and implementation, and A/B testing the efficacy of intervention parameters before deploying an experiment at large. That said, these tradeoffs permit an innovative and powerful approach for causal investigations of attention’s impacts on subsequent cognitive behavior, advancing understanding of the multiple processes that collectively influence whether and how we learn and remember experiences.

\section{Concluding Remarks}

This article highlighted some recent advances in understanding the neurocognitive mechanisms of attention, episodic memory, and their interactions. We expect that future research and further methodological developments will continue to drive theoretical progress, enabling researchers to tackle the open questions raised herein and generate increasingly precise accounts of the systems and mechanisms that enable humans to learn and remember from life events, as well as how mnemonic function changes in aging and disease.

% \section{Recommended Reading}

% \begin{itemize}
%     \item Kerrén, C., Linde-Domingo, J., Hanslmayr, S., \& Wimber, M. (2018). (See References). Scalp EEG study demonstrating fluctuations in memory reinstatement evidence within the theta frequency band.
%     \item Korkki, S. M., Richter, F. R., Jeyarathnarajah, P., \& Simons, J. S. (2020). (See References). Multiexperiment study demonstrating age-related reduction of episodic memory precision across a number of dimensions.
%     \item Madore, K. P., \& Wagner, A. D. (2022). (See References). Detailed theoretical account of how variability in sustained attention within and across individuals influences one’s readiness to learn and remember.
%     \item VanRullen, R. (2018). Attention cycles. \textit{Neuron, 99}(4), 632–634. Concise summary of a number of nonhuman and human studies linking electrophysiological oscillations with behavioral rhythms of attention.
%     \item Xia, R., Chen, X., Engel, T. A., \& Moore, T. (2024). Common and distinct neural mechanisms of attention. \textit{Trends in Cognitive Sciences, 28}(6), 554–567. Highlights advances in the understanding of stimulus-driven and top-down attentional control.
% \end{itemize}

% \section{Acknowledgments}
% We acknowledge the contributions of the many authors whose work we were unable to include as a result of space and citation limitations. We also extend our gratitude to the past and present members of the Stanford Memory Lab who have influenced the work discussed herein. Funding This work was supported by an agility grant from the Stanford University Wu Tsai Human Performance Alliance (to A. D. Wagner \& S. T. Schwartz); the Stanford University Wu Tsai Neurosciences Institute Center for Mind, Brain, Computation and Technology (to S. T. Schwartz \& A. M. Xue); a Stanford University Ric Weiland Graduate Fellowship in the Humanities \& Sciences (to S. T. Schwartz); the McKnight Brain Research Foundation Clinical Translational Research Scholarship in Cognitive Aging and Age-Related Memory Loss; the American Brain Foundation; the American Academy of Neurology (to H. Yang); National Institutes of Health Grant R01-AG065255 (to A. D. Wagner); and National Science Foundation Graduate Research Fellowship DGE-2146755 (to A. M. Xue).

